% TOP_PROBLEM: {"problemId": "problem.equivariant-tamagawa-number-conjecture", "canonicalTitle": "Equivariant Tamagawa number conjecture: Burns-Flach Conjecture4 package", "releaseRank": 58, "primaryDomain": "number_theory_arithmetic_geometry", "primaryDomainLabel": "Number theory & arithmetic geometry", "displayStatus": "open_with_solved_subcases", "releaseStatus": "open_with_solved_subcases"}
% !TeX program = lualatex
% Definition and sources only; this file is not an LLM solution attempt.
\documentclass[11pt]{article}
\usepackage[margin=25mm]{geometry}
\usepackage{fontspec}
\setmainfont{Latin Modern Roman}
\usepackage{amsmath,amssymb,mathtools}
\usepackage{unicode-math}
\setmathfont{Latin Modern Math}
\usepackage{xurl,hyperref}
\hypersetup{colorlinks=true,urlcolor=blue}
\setcounter{secnumdepth}{0}
\setlength{\parindent}{0pt}
\setlength{\parskip}{5pt}
\setlength{\emergencystretch}{3em}
\begin{document}
\section*{\texorpdfstring{Equivariant Tamagawa number conjecture: Burns-Flach Conjecture4 package}{Equivariant Tamagawa number conjecture: Burns-Flach Conjecture4 package}}\label{58}
\subsection{Definitions and mathematical statement}
Let $A$ be a finite-dimensional semisimple ${\mathbb{Q}}$-algebra and $\mathcal A\subset A$ an order with the compatible projective structure on $M$ specified by Burns--Flach. Relative $K_0(\mathcal A,{\mathbb{R}})$ records projective modules together with isomorphisms after extension to $A\otimes{\mathbb{R}}$. The conjecture is the \emph{entire} conditional package of Conjecture 4(i)--(iv): meromorphic continuation at zero, the reduced-rank order formula, rationality of the Tamagawa class, and its vanishing. In particular,
\[
 r=\operatorname{rr}_A H_f^1(K,M^\vee(1))^*
   -\operatorname{rr}_A H_f^0(K,M^\vee(1))^*,
\]
\[
 T_\Omega=\widehat\delta^1_{\mathcal A,{\mathbb{R}}}\bigl(L^*(AM,0)\bigr)
             +R_\Omega(M,\mathcal A)=0.
\]
Here $L^*$ is the componentwise leading Laurent coefficient, $R_\Omega$ is the fundamental virtual-object class with its comparison trivialization, and $\widehat\delta$ is the canonical extended reduced-norm boundary. Their exact definitions and the external comparison, finiteness, coefficient-compatibility, and coherence hypotheses are incorporated from the source. Noncommutative and nonmaximal orders remain in scope. The full recorded target below prevents this last identity from being mistaken for the whole package.
\par\smallskip\textit{Formulation and concept references:} \href{https://ems.press/content/serial-article-files/25892}{[S1]}.

\subsection{Short English statement}
Do equivariant motivic $L$-functions satisfy the full Burns--Flach continuation, vanishing-order, rationality, and integral leading-term law under its stated prerequisites?
\subsection{Sources}
\begin{itemize}
\item[S1]David Burns and Matthias Flach, Tamagawa Numbers for Motives with (Non-Commutative) Coefficients, Documenta Mathematica6(2001),501-570.
\url{https://ems.press/content/serial-article-files/25892}.
\textit{Formulation source. Consulted 24 September 2026: Sections 3--4; relative K-theory and extended boundary.}
\item[S2]The Tamagawa number conjecture and Kolyvagin's conjecture for motives of modular forms.
\url{https://link.springer.com/article/10.1007/s40687-026-00646-7}.
\textit{Further reference; not a proof endorsement.}
\item[S3]An unconditional proof of the abelian equivariant Iwasawa main conjecture and applications.
\url{https://athene-forschung.unibw.de/92554?show_id=155671}.
\textit{Further reference; not a proof endorsement.}
\end{itemize}
\end{document}
